\documentclass[11pt]{article}

\usepackage[utf8]{inputenc}
\usepackage[T1]{fontenc}
\usepackage{lmodern}
\usepackage[margin=1in]{geometry}
\usepackage{amsmath,amssymb,amsthm}
\usepackage{microtype}
\usepackage{booktabs}
\usepackage{xcolor}
\usepackage{listings}
\usepackage{enumitem}
\usepackage[hidelinks]{hyperref}
\usepackage{cleveref}

\newtheorem{theorem}{Theorem}
\newtheorem{proposition}[theorem]{Proposition}
\newtheorem{corollary}[theorem]{Corollary}
\theoremstyle{definition}
\newtheorem{definition}[theorem]{Definition}
\theoremstyle{remark}
\newtheorem{remark}[theorem]{Remark}

\newcommand{\R}{\mathbb{R}}
\newcommand{\Q}{\mathbb{Q}}
\newcommand{\F}{\mathbb{F}}
\newcommand{\RD}{\mathrm{RD}}
\newcommand{\RU}{\mathrm{RU}}
\newcommand{\lean}[1]{\texttt{#1}}
\newcommand{\realfp}{\textsf{realfp}}

\definecolor{kw}{RGB}{0,0,160}
\definecolor{cm}{RGB}{100,100,100}
\lstdefinelanguage{lean}{
  morekeywords={theorem,example,by,import,namespace,end,fun,def,open},
  sensitive=true,
  morecomment=[l]{--},
  morecomment=[s]{/-}{-/},
  morestring=[b]",
}
\lstset{
  language=lean,
  basicstyle=\ttfamily\small,
  keywordstyle=\color{kw}\bfseries,
  commentstyle=\color{cm}\itshape,
  columns=fullflexible,
  keepspaces=true,
  breaklines=true,
  xleftmargin=1em,
  aboveskip=0.6em, belowskip=0.6em,
  extendedchars=true,
  literate=
    {∀}{{$\forall$}}1 {∃}{{$\exists$}}1 {∈}{{$\in$}}1 {ℝ}{{$\mathbb{R}$}}1
    {ℚ}{{$\mathbb{Q}$}}1 {ℕ}{{$\mathbb{N}$}}1 {≤}{{$\le$}}1 {≥}{{$\ge$}}1
    {∫}{{$\textstyle\int$}}1 {∧}{{$\wedge$}}1 {×}{{$\times$}}1 {→}{{$\to$}}1
    {≠}{{$\ne$}}1 {ˢ}{{$^{s}$}}1 {π}{{$\pi$}}1 {√}{{$\surd$}}1
    {²}{{$^2$}}1 {⁻}{{$^{-}$}}1 {¹}{{$^1$}}1,
}

\title{\realfp: Fast, Kernel-Checked Interval Arithmetic\\ for Lean~4 and Mathlib}
\author{Roger Schlafly\thanks{E-mail: \href{mailto:roger@darkbuzz.com}{\texttt{roger@darkbuzz.com}}.
ORCID: \href{https://orcid.org/0009-0006-9524-245X}{\texttt{0009-0006-9524-245X}}.}}
\date{November 2026}

\emergencystretch=2em
\begin{document}
\maketitle

\begin{abstract}
\realfp{} is an open-source Lean~4 library for proving inequalities, root
existence and uniqueness, integral bounds, positive definiteness and extremal
values of real-valued expressions by interval arithmetic. Interval endpoints
are IEEE~754 binary64 numbers; every directed-rounding operation is a
computable Lean definition proved correct against Mathlib's real numbers, and
compiled code runs the same operations through a small C runtime that uses
hardware directed rounding (AVX-512 embedded rounding on x86-64, or a portable
round-to-nearest plus exact-error-term method elsewhere). A proof tactic
splits each goal into an untrusted native search for a certificate and a
verified Boolean checker. By default the checker is evaluated by Lean's
kernel, so the resulting theorems depend only on Lean's three standard axioms
and do not trust the C code, the compiler or the processor. Kernel checking
uses natural-number mirrors of the floating-point operations that are proved
equal to the specification. An opt-in native mode trades a larger trusted
base for speed on very large certificates, and a 128-bit dyadic mode gives
about 38 significant digits. Beyond plain bisection, \realfp{} uses derivative
and monotonicity certificates, Krawczyk's operator, Taylor-model quadrature,
exact rational $LDL^{\mathsf T}$ factorizations, and a second-order
certificate for bounds that are attained at interior points. We describe the
design, state the soundness theorems, give examples, compare with Arb/FLINT,
CoqInterval, FloatLib and LeanCert, and report benchmarks.
\end{abstract}

\section{Introduction}

Computer-assisted proofs in analysis often come down to a large number of
explicit numerical facts: a function is positive on a box, a system has
exactly one zero in a region, an integral lies in a given interval.
Interval arithmetic~\cite{moore2009} turns such facts into theorems: if every
operation returns an interval guaranteed to contain the exact result, an
enclosure that excludes zero proves that the quantity is not zero. Rigorous
numerics of this kind was essential to Tucker's proof that the Lorenz
attractor exists~\cite{tucker2002} and to the Kepler conjecture, whose
formal proof (Flyspeck) verified hundreds of nonlinear inequalities in HOL
Light~\cite{flyspeck,solovyevhales2013}.

Lean~4~\cite{lean4} with Mathlib~\cite{mathlib} has become a common system
for formalizing mathematics, but its native \lean{Float} type is opaque to
the kernel, so nothing can be proved about a hardware floating-point
computation. \realfp{}~\cite{realfp} closes that gap. It is meant for the
mathematician who has a numerical lemma inside a larger formal proof and
wants it proved quickly, with a small and clearly stated trusted base, by a
library that can simply be added as a Lake dependency. Two lines suffice for
typical goals:
\begin{lstlisting}
example : ∀ x ∈ Set.Icc (0 : ℝ) 1, Real.sin x ≤ x + 1/1000 := by realfp
example : ∃! x ∈ Set.Icc (1 : ℝ) 2, x ^ 2 = 2 := by realfp
\end{lstlisting}
Both theorems are checked by the Lean kernel and depend only on
\lean{propext}, \lean{Classical.choice} and \lean{Quot.sound}.

\paragraph{Contributions.}
\begin{itemize}[nosep]
\item A binary64 interval library whose directed roundings are Lean
  definitions with machine-checked enclosure theorems, and whose compiled
  implementation uses hardware directed rounding (\cref{sec:design}).
\item A kernel fast path: natural-number mirrors of all operations, proved
  equal to the specification, so that certificates are checked by
  \lean{decide +kernel} without trusting compiled code (\cref{sec:kernel}).
\item One front-door tactic, \lean{realfp}, over a family of verified
  certificate checkers: bisection, derivative and monotonicity certificates,
  multivariate centred forms, Krawczyk systems, Taylor-model integration,
  exact positive-definiteness, optimization, and touching-point
  certificates; and a 128-bit dyadic mode for all of them (\cref{sec:tactics}).
\item An assurance story aimed at users: no \lean{sorry}, a build-time axiom
  audit of every theorem, bit-for-bit differential tests of the C runtime,
  and continuous integration (\cref{sec:assurance}).
\item Benchmarks against other Lean tools and a downstream re-verification
  of a published formal proof (\cref{sec:bench}).
\end{itemize}
This paper describes version~0.1.1, which adds the touching-point
features and build portability to the first release~0.1.0. The library is licensed under
Apache-2.0; full documentation, benchmarks scripts and logs are in the
repository~\cite{realfp}.

\section{Background}

\paragraph{Directed rounding.}
Let $\F$ be the set of binary64 numbers, extended by $\pm\infty$. IEEE~754
\cite{ieee754} specifies, for each basic operation $\circ\in\{+,-,\times,\div\}$
and for $\sqrt{\cdot}$, that the result is the exact real result rounded by a
stated rule. The rules $\RD$ (toward $-\infty$) and $\RU$ (toward
$+\infty$) give $\RD(a\circ b)\le a\circ b\le\RU(a\circ b)$. An interval
$[\ell,u]$ with $\ell,u\in\F$ denotes $\{x\in\R:\ell\le x\le u\}$, and
outward rounding yields the inclusion property $x\in X,\ y\in Y\Rightarrow
x\circ y\in X\circledcirc Y$. Composing such operations along an expression
gives an enclosure of its range over a box; subdividing the box tightens it.
On x86-64 the AVX-512 instruction set allows the rounding direction to be
encoded in each instruction (``embedded rounding''), so no global rounding
mode has to be switched.

\paragraph{Kernel and native evaluation in Lean.}
A Lean proof of a decidable proposition $P$ can be produced by evaluation.
\lean{decide +kernel} asks the kernel to reduce the decision procedure to
\lean{true}; the result depends on nothing but the kernel and the standard
axioms. Lean's kernel has GMP-accelerated natural-number arithmetic, which
makes evaluation on \lean{Nat} encodings fast. \lean{native\_decide} instead
compiles the decision procedure, runs it, and admits the result through an
auxiliary axiom (generated from \lean{Lean.ofReduceBool}); this trusts the
Lean compiler, the C compiler, and any code attached with
\lean{@[implemented\_by]} or \lean{@[extern]}. \realfp{} supports both,
with kernel evaluation as the default.

\section{Design}\label{sec:design}

\subsection{Specification in Lean, speed in C}

Every directed operation (\lean{addDown}, \lean{mulUp}, \lean{sqrtDown},
\dots) is a plain, computable Lean definition: a software implementation of
IEEE binary64 rounding toward $\pm\infty$ on the bit patterns, built on
Lean's \lean{Float} model. A C kernel is attached to each definition with
\lean{@[implemented\_by]}. Compiled code (\lean{\#eval}, the certificate
search inside the tactics, and \lean{native\_decide}) runs the C version,
while the kernel sees only the Lean definition. There are two C backends:
\begin{itemize}[nosep]
\item \emph{AVX-512.} Each operation is a single instruction with embedded
  $\RD$ or $\RU$ rounding (e.g.\ \lean{\_mm\_add\_round\_sd}).
\item \emph{Portable RN+error.} The result $s=\mathrm{RN}(a\circ b)$ is
  computed in the default round-to-nearest mode together with its exact
  error term: TwoSum for addition, $\mathrm{fma}(a,b,-s)$ for
  multiplication, the residual $\mathrm{fma}(-q,b,a)$ for division and
  $\mathrm{fma}(-s,s,x)$ for square root. The sign of the error decides
  whether to step to the adjacent float. Rare cases where the error term
  itself could underflow (tiny operands) take a scaled or exact slow path.
  This backend is the default on arm64 and is used in CI.
\end{itemize}
Both backends are required to agree bit for bit with the Lean
specification, and they are tested to do so (\cref{sec:assurance}). The
theorems never depend on which backend is used.

Elementary functions ($\exp,\log,\log_2,\log_{10},\sin,\cos,\tan,\arctan,
\arcsin,\arccos,\sinh,\cosh,\tanh$, $\pi$, real powers) have a Lean
specification given by rational enclosures with explicit error terms,
refined by a Ziv loop until the directed rounding of both ends agrees. At
run time they use the correctly rounded CORE-MATH functions~\cite{coremath}
executed under $\RD$ or $\RU$ hardware rounding.

\subsection{Interval semantics}

An \lean{Interval64} is a pair of binary64 bit patterns. Its meaning is a set
of reals,
\[
  \lean{X.toSet}=\{x\in\R:\ \lean{lo}\le x\le\lean{hi}\},
\]
where an endpoint is read as an extended real ($-\infty$ and $+\infty$
are allowed, so overflow is harmless) and a NaN endpoint makes the set
empty. Expressions are reified into an inductive type \lean{IExpr} with
variables, rational constants, $+,-,\times,\div$, negation, inverse,
$\sqrt{\cdot}$, $|\cdot|$, natural powers, the elementary functions above,
$\pi$, $\operatorname{atan2}$, and user-registered unary functions. Its
denotation \lean{e.denote xs} uses Mathlib's (total) real functions, so for
example $\log 0=0$ and $1/0=0$ are covered by the enclosures as Mathlib
defines them.

\subsection{The kernel fast path}\label{sec:kernel}

Reducing the bit-level specification in the kernel is correct but slow
(about 3\,ms per operation, \cref{tab:ops}). \realfp{} therefore defines
\lean{KF} and \lean{KI}: mirrors of the scalar and interval operations
written directly on \lean{Nat}, which the kernel evaluates with GMP. They
are proved equal to the specification
(\lean{Interval64Proofs/KernelEq.lean} and neighbours); for expressions,
\[
  \lean{toI}\,(\lean{evalK}\ \mathit{env}\ e) \;=\; \lean{evalI}\,(\mathit{env}.\lean{map toI})\ e ,
\]
where \lean{evalI} is the specification-level evaluator
(\lean{evalK\_toI}). Speed in the kernel therefore costs nothing in trust.

\subsection{Search and check}

Every tactic splits its work in two.
\emph{Search} (untrusted) runs compiled code with the fast C kernels and
produces a certificate: a bisection tree, derivative signs, a Krawczyk box
and preconditioner, a quadrature partition, an $LDL^{\mathsf T}$ factor.
A bug in the search can only make the tactic fail.
\emph{Check} (trusted) is a Boolean function of the goal data and the
certificate, proved sound in Lean/Mathlib, and closed by
\lean{decide +kernel} (default) or \lean{native\_decide} (opt-in).

\subsection{128-bit mode}

When 16 digits are not enough, \lean{(precision := 128)} switches every
tactic to 128-bit dyadic intervals (\lean{Interval128.Dy}): endpoints
$m\cdot 2^{e}$ with 128-bit mantissas, implemented directly on \lean{Nat}
and \lean{Int} for the kernel, so there is no separate C runtime to trust.
Transcendental functions $\exp,\log,\log_2,\log_{10},\sin,\cos,\sinh,\cosh,
\tanh$ are evaluated at this precision; $\pi$ comes from 40-digit rational
bounds proved from Mathlib's. Two implementation details matter for kernel
speed: split points are emitted as sums of 62-bit limbs (the kernel is slow
on literals $\ge 2^{63}$), and 128-bit integration uses fixed-point Taylor
models whose products are each certified by a witness, so the kernel never
re-evaluates shared intermediate terms.

\section{Soundness theorems}\label{sec:theorems}

We state the main theorems in mathematical notation and give their Lean
names. All of them are proved in Lean without \lean{sorry} and depend only
on the three standard axioms.

\begin{theorem}[Directed rounding; \lean{add\_encloses}, \lean{div\_encloses}, \dots]
\label{thm:round}
Let $a,b$ be finite binary64 numbers with rational values $q_a,q_b$, and
let $\circ\in\{+,-,\times,\div\}$ (with $q_b\ne0$ for $\div$). Then
\[
  \lean{opDown}(a,b)\ \le\ q_a\circ q_b\ \le\ \lean{opUp}(a,b)
\]
as extended reals, and both sides equal $q_a\circ q_b$ when it is
representable. The analogous statements hold for $\sqrt{\cdot}$.
\end{theorem}
These are proved in the core package \lean{Interval64Proofs}, which does not
depend on Mathlib, by reasoning about the bit-level specification.

\begin{theorem}[Inclusion; \lean{add\_mem\_toSet}, \dots, \lean{eval\_sound}]
\label{thm:eval}
For every expression $e$ and lists $x=(x_1,\dots,x_n)$ of reals and
$X=(X_1,\dots,X_n)$ of intervals with $x_i\in X_i$,
\[
  e.\lean{denote}(x)\ \in\ \lean{evalI}(X,e).\lean{toSet}.
\]
\end{theorem}
The proof is by induction on $e$; each case is an inclusion lemma for one
operation over $\R$, such as $x\in X,\ y\in Y\Rightarrow x+y\in X+Y$, with
infinite endpoints allowed.

\begin{theorem}[Certificate soundness; \lean{prove\_le}, \lean{checkK\_sound}]
\label{thm:cert}
Let $e$ be an expression, $B=\prod_i[a_i,b_i]$ a box with rational
endpoints, $c\in\Q$ and $\mathcal C$ a bisection certificate. If
$\lean{checkK}(e,\ \lean{gLe}\ c,\ \mathcal C,\ B)=\lean{true}$, then
$e.\lean{denote}(x)\le c$ for every $x\in B$. The same holds for $<$, $\ge$,
$>$ and $\ne$ (\lean{prove\_lt}, \lean{prove\_ge}, \lean{prove\_gt},
\lean{prove\_ne}).
\end{theorem}

In kernel mode the hypothesis $\lean{checkK}(\dots)=\lean{true}$ is
discharged by \lean{decide +kernel}. Combining
\cref{thm:round,thm:eval,thm:cert} with the kernel mirror gives:

\begin{corollary}[Trusted base]\label{cor:trust}
A theorem proved by a \realfp{} tactic in the default mode
\lean{trust := kernel} depends only on \lean{propext},
\lean{Classical.choice} and \lean{Quot.sound}. It trusts the Lean kernel
(including its GMP arithmetic) and the meaning of the Mathlib definitions in
its statement; it does not trust the C runtime, any compiler, the
floating-point hardware, or the search code. In \lean{trust := native}
mode the theorem additionally depends on generated
\lean{native\_decide} axioms (one per checked chunk), and the proof additionally trusts the
Lean and C compilers, the agreement of each C kernel with its Lean
specification, CORE-MATH's correct rounding under directed modes, and the
default floating-point environment (no flush-to-zero).
\end{corollary}
The corollary is not a Lean theorem but a statement about what
\lean{\#print axioms} reports; it is enforced mechanically by the axiom audit
described in \cref{sec:assurance}.

The other certificate families have soundness theorems of the same shape.
We list the mathematical content.

\begin{theorem}[Derivative certificates; \lean{checkD\_sound}, \lean{checkM\_sound}]
If, on a sub-box, the interval enclosure of $\partial f/\partial x_k$ has
constant sign, the bound may be checked on the corresponding face; if
$f(m)+\nabla f(X)\cdot(X-m)$ (a centred form) satisfies the bound, so does
$f$ on $X$. Differentiability on the box is checked by a verified domain
predicate (\lean{domOK\_sound}), and derivatives are computed symbolically
by \lean{IExpr.deriv}, proved correct against Mathlib's \lean{HasDerivAt}.
\end{theorem}

\begin{theorem}[Krawczyk; \lean{kraw\_exists}, \lean{kraw\_unique}]
Let $F:\R^n\to\R^n$ be given by expressions, $X$ a rational box, $m\in X$,
$Y,Z$ rational matrices and $q<1$. If the checker verifies
$m-YF(m)+(I-YJ)(X-m)\subseteq X$, the row sums of $|I-YJ|$ are at most $q$,
and $\|I-ZY\|_\infty\le q$, where $J$ is an interval enclosure of the
Jacobian on $X$, then $F$ has exactly one zero in $X$.
\end{theorem}
The proof applies Banach's fixed-point theorem to $G(x)=x-YF(x)$ on $X$;
injectivity of $Y$ turns the fixed point into a zero of $F$~\cite{krawczyk1969}.

\begin{theorem}[Integration; \lean{integral\_bounds\_of\_quadBound}, \lean{integral\_bounds\_of\_quadBoundF}]
If a partition $a=t_0<\dots<t_N=b$ and, on each cell, a Taylor
expansion of order $K$ about the midpoint with an interval remainder from
$f^{(K)}$ on the cell pass the checker, then
$\ell\le\int_a^b f(x)\,dx\le h$. In 128-bit mode the Taylor models are
fixed-point polynomials with a rigorous radius, and each product carries
a witness that the kernel verifies~\cite{makinoberz2003}.
\end{theorem}

\paragraph{Touching points.}
Interval bisection cannot prove $f\ge c$ if $f(p)=c$ at an interior point
$p$: every box containing $p$ has an enclosure whose lower end is slightly
below $c$. \realfp{} handles such goals with a second-order certificate.

\begin{theorem}[Tangent-plane leaf; \lean{touch\_expand\_gen}, \lean{quad\_nonneg\_of\_cert}, \lean{tanOK\_sound}]
\label{thm:touch}
Let $B$ be a box containing $p$, and let $H$ be the interval enclosure of
the Hessian of $f$ on $B$. Suppose there are a rational symmetric matrix $C$
and rationals $\rho_k$ such that $\sum_j|S_{kj}-C_{kj}|\le\rho_k$ for every
real symmetric $S$ with entries in $H$, and $C-\operatorname{diag}(\rho)$
is positive definite (checked exactly over $\Q$ by an $LDL^{\mathsf T}$
congruence). Then for all $x\in B$,
\[
  f(x)\ \ge\ f(p)+\nabla f(p)\cdot(x-p).
\]
Hence $f\ge c$ on $B$ follows from $f(p)+\nabla f(p)\cdot(x-p)\ge c$ on $B$,
which is checked by interval arithmetic.
\end{theorem}
The proof expands $f$ coordinate by coordinate with the one-variable mean
value theorem, producing $2(f(x)-f(p))=2\nabla f(p)\cdot u+u^{\mathsf T}Su$
with $u=x-p$ and $S$ symmetric with entries in $H$; the certificate makes
$u^{\mathsf T}Su\ge0$. When $p$ is rational and $f$ uses only
$+,-,\times,\div,|\cdot|$ and powers, $f(p)$ and $\nabla f(p)$ are
evaluated exactly in $\Q$ (\lean{evalQ\_sound}), so the tangent plane at
an exact touching point such as $(1/3,1/5)$ is exact. With
\lean{touch := auto} the points are found by an untrusted multistart Newton
search on $\nabla f=0$.

\section{Tactics}\label{sec:tactics}

\lean{realfp} inspects the goal and dispatches:
\begin{itemize}[nosep]
\item $\forall x\in B,\ f\bowtie g$ with ${\bowtie}\in\{\le,<,\ge,>,\ne\}$, and
  closed inequalities between constants: \lean{interval} (bisection, then
  derivative, monotonicity and centred-form certificates, Horner forms for
  polynomials);
\item $\exists x\in[a,b],\ f(x)=0$ and $\exists!$: \lean{root\_exists},
  \lean{root\_unique} (sign change plus monotonicity in one variable,
  Newton search plus Krawczyk for square systems);
\item $\int_a^b f\in[\ell,h]$: \lean{integral\_bound} (adaptive Taylor
  quadrature);
\item \lean{Matrix.PosDef}: \lean{pos\_def} (interval or exact rational
  $LDL^{\mathsf T}$ congruence certificates);
\item $\operatorname{sInf}/\operatorname{sSup}$ of $f(B)$, existence of
  $\exists m\,\forall x$ bounds, witnesses and one-variable argmin:
  branch-and-bound search, then the same certificates.
\end{itemize}
Options include \lean{trust := kernel | native}, \lean{precision := 64 | 128},
\lean{depth}, and \lean{touch := [p$_1$, \dots] | auto}. Failure messages name
the failing box in the user's variables and report whether the bound appears
false (with an estimate of the true minimum) or merely too tight.
\lean{\#realfp\_min} and \lean{\#realfp\_max} report numeric extrema
interactively. New unary functions can be registered with an attribute by
supplying an interval extension, its derivative and a soundness proof; the
library itself registers $\operatorname{erf}$, $\operatorname{arsinh}$,
$\operatorname{artanh}$ and $\operatorname{sinc}$ this way.

\section{Examples}\label{sec:examples}

All snippets in this section are compiled in CI and, unless noted, are
checked by the kernel with the three standard axioms; the times are file
elaboration times measured for \cref{tab:kn}.

\paragraph{Constants.} Thirty digits of $\pi$ need the 128-bit mode;
Archimedes' bound and $e^\pi>\pi^e$ do not.
\begin{lstlisting}
theorem pi_bounds :
    (3.14159265358979323846264338327950 : ℝ) < Real.pi ∧
      Real.pi < 3.14159265358979323846264338327951 := by
  realfp (precision := 128)

theorem pi_lt_22_div_7 : Real.pi < 22 / 7 := by realfp

theorem exp_pi_gt_pi_exp :
    Real.exp (Real.exp 1 * Real.log Real.pi) < Real.exp Real.pi := by realfp
\end{lstlisting}

\paragraph{The Gaussian integral.} $\int_0^1 e^{-x^2}dx=0.746824132812427025\ldots$
has no elementary antiderivative.
\begin{lstlisting}
theorem gauss_four_digits :
    (∫ x in (0 : ℝ)..1, Real.exp (-x ^ 2)) ∈ Set.Icc (0.7468 : ℝ) 0.7469 := by
  integral_bound

theorem gauss_18_digits :
    (0.746824132812427025 : ℝ) ≤ ∫ x in (0 : ℝ)..1, Real.exp (-x ^ 2) ∧
      (∫ x in (0 : ℝ)..1, Real.exp (-x ^ 2)) ≤ 0.746824132812427026 := by
  integral_bound (precision := 128)
\end{lstlisting}

\paragraph{Jordan's inequality.} The tactic needs numeral endpoints, so the
interval is $[0,1.57]\subset[0,\pi/2]$.
\begin{lstlisting}
theorem jordan :
    ∀ x ∈ Set.Icc (0 : ℝ) (157 / 100), 2 * x / Real.pi ≤ Real.sin x := by
  realfp
\end{lstlisting}

\paragraph{Touching points.} AM--GM, $1+x\le e^x$ and the Motzkin
polynomial~\cite{motzkin1967} $M(x,y)=x^4y^2+x^2y^4-3x^2y^2+1$, which is
nonnegative but not a sum of squares, all attain their bound at interior
points. On $[-2,2]^2$, $M\ge0$ follows from the quadrant by symmetry; with
\lean{touch := auto} the four zeros $(\pm1,\pm1)$ are also found
automatically.
\begin{lstlisting}
theorem sqrt_le_avg :
    ∀ x ∈ Set.Icc (0 : ℝ) 10, Real.sqrt x ≤ (1 + x) / 2 := by
  realfp (touch := [1])

theorem one_add_le_exp : ∀ x ∈ Set.Icc (-2 : ℝ) 5, 1 + x ≤ Real.exp x := by
  realfp (touch := auto)

theorem motzkin_quadrant :
    ∀ x ∈ Set.Icc (0 : ℝ) 2, ∀ y ∈ Set.Icc (0 : ℝ) 2,
      0 ≤ x ^ 4 * y ^ 2 + x ^ 2 * y ^ 4 - 3 * x ^ 2 * y ^ 2 + 1 := by
  realfp (touch := [(1, 1)])
\end{lstlisting}
The certificate also covers \emph{near} touching at irrational points. The
polynomial below is $(x^2-2)^2+(3y^2-1)^2+10^{-13}$, whose minimum is at
$(\sqrt2,1/\sqrt3)$; plain bisection fails, and \lean{touch := auto} proves
it with about 40 boxes.
\begin{lstlisting}
theorem near_touch :
    ∀ x ∈ Set.Icc (0 : ℝ) 2, ∀ y ∈ Set.Icc (0 : ℝ) 2,
      0 < x ^ 4 - 4 * x ^ 2 + 4 + 9 * y ^ 4 - 6 * y ^ 2 + 1 + 1 / 10 ^ 13 := by
  realfp (touch := auto)
\end{lstlisting}

\paragraph{Roots, systems, extrema, matrices.}
\begin{lstlisting}
theorem dottie_unique : ∃! x ∈ Set.Icc (0 : ℝ) 1, Real.cos x = x := by realfp

theorem sys2_unique : ∃! p : ℝ × ℝ, p.1 ∈ Set.Icc (0.9 : ℝ) 1.1 ∧
    p.2 ∈ Set.Icc (0.9 : ℝ) 1.1 ∧ p.1 ^ 2 + p.2 = 2 ∧ p.1 + p.2 ^ 2 = 2 := by
  realfp

theorem exp_sInf :
    sInf ((fun x : ℝ => Real.exp x - 2 * x) '' Set.Icc 0 2) ∈
      Set.Icc 0.613705 0.613706 := by
  realfp

example : (!![4, 1, 0; 1, 3, 1; 0, 1, 2] : Matrix (Fin 3) (Fin 3) ℝ).PosDef := by
  pos_def
\end{lstlisting}

\paragraph{Native mode.} With \lean{(trust := native)}, \lean{\#print axioms}
lists the three standard axioms and one generated axiom,
\lean{exp\_pi\_native.\_native.\allowbreak native\_decide.ax\_1\_7}:
\begin{lstlisting}
theorem exp_pi_native :
    Real.exp 1 * Real.log Real.pi < Real.pi := by realfp (trust := native)
\end{lstlisting}

\section{Benchmarks}\label{sec:bench}

Unless stated otherwise, measurements were taken on 25~November~2026 on a
shared Linux virtual machine with 8 vCPUs (Intel Xeon with AVX-512F) and
15\,GiB of memory, with Lean and Mathlib \lean{v4.36.0}, gcc~14.2 and the
AVX-512 backend, at \realfp{} commit \lean{61f8eef}. The scripts, raw logs
and further comparisons are in the repository; numbers taken from earlier
logged runs are marked as such.

\paragraph{Kernel versus native checking.}
\Cref{tab:kn} gives the elaboration time of single theorems, measured as
the wall time of \lean{lake lean} on a file containing only that theorem
(median of three runs) minus the time to load the imports (3.7\,s). The
times include the untrusted search, which is the same in both modes. Every
kernel-mode theorem depends only on the three standard axioms; every
native-mode theorem depends in addition on one or more generated
\lean{native\_decide} axioms (one per checked chunk, e.g.\ 26 for the camel
function).

\begin{table}[ht]
\centering\small
\caption{Elaboration time of single theorems, in seconds, above the 3.7\,s
import baseline (``$<0.5$'': indistinguishable from the baseline).}
\label{tab:kn}
\begin{tabular}{@{}llrr@{}}
\toprule
Goal & Kind & kernel & native\\
\midrule
$\pi<22/7$ & constant & $<0.5$ & $<0.5$\\
$\pi$ to 30 digits & constant, 128-bit & $<0.5$ & $<0.5$\\
$e^\pi>\pi^e$ & constant & 2.1 & $<0.5$\\
$\exists!$ zero of a $2\times2$ system & root of a system & $<0.5$ & $<0.5$\\
$e^{xy}-x-y\ge-0.01$ on $[0,1]^2$ & 2-variable bound & 4.0 & 2.2\\
six-hump camel $\ge-1.032$ on $[-2,2]\times[-1,1]$ & 2-variable polynomial & 22.2 & 3.6\\
Jordan, $[0,1.57]$ & 1-variable bound & 17.4 & $<0.5$\\
$\int_0^1e^{-x^2}$, 4 digits & integral & 3.6 & 1.2\\
$\int_0^1e^{-x^2}$, 18 digits & integral, 128-bit & 4.7 & $<0.5$\\
Motzkin $\ge0$ on $[0,2]^2$ & touching point & 15.9 & 2.5\\
\bottomrule
\end{tabular}
\end{table}

For small and medium certificates kernel checking costs a few seconds,
which is why it is the default. The gap grows with the number of
transcendental evaluations: Jordan's inequality needs many enclosures of
$\sin$ on small boxes, each of which the kernel evaluates through rational
series.

\paragraph{Cost of one operation.}
\Cref{tab:ops} shows the cost of a single directed operation in compiled
code and in the kernel. Compiled times are from the repository's
micro-benchmark (a dependent chain of $2\cdot10^7$ upward additions, one
core); the elementary functions are CORE-MATH under directed rounding,
called from Lean. Kernel times use FloatLib's kernel-reduction harness
(\lean{benchmarks/kernel}, \lean{Lean.Kernel.whnf} on a table of eight
binary64 operand triples with normal operands, fastest of three, a no-op
over the same table subtracted), run on the same machine for both
libraries: \realfp's \lean{KF} mirror rounding downward, and FloatLib's
kernel path in its \lean{down} direction (revision \lean{d24844f}).

\begin{table}[ht]
\centering\small
\caption{Cost of one binary64 directed operation.}
\label{tab:ops}
\begin{tabular}{@{}lrrrrr@{}}
\toprule
 & add & sub & mul & div & sqrt\\
\midrule
\multicolumn{6}{@{}l}{\emph{Compiled code (ns)}}\\
\quad \realfp{} AVX-512 embedded rounding & 1.2 & & & & \\
\quad \realfp{} portable RN+error & 5.3 & & & & \\
\quad \lean{fesetround} around each operation (reference) & 14.3 & & & & \\
\midrule
\multicolumn{6}{@{}l}{\emph{Lean kernel ($\mu$s)}}\\
\quad \realfp{} \lean{KF} (\lean{Nat} mirror) & 189 & 191 & 182 & 227 & 240\\
\quad FloatLib kernel path & 87 & 73 & 51 & 61 & 52\\
\quad \realfp{} bit-level specification$^{*}$ & 4900 & 4400 & 4800 & 3700 & 3800\\
\bottomrule
\end{tabular}

\smallskip
\parbox{0.92\linewidth}{\footnotesize $^{*}$Single \lean{Kernel.isDefEq}
check of one result, median of ten operand pairs (a different harness);
shown only to explain why the mirror exists. CORE-MATH $\exp$, $\log$,
$\sin$ under directed rounding cost 17--21\,ns per call in compiled code.}
\end{table}

Two points stand out. Hardware directed rounding is about $10^5$ times
faster than kernel evaluation, which is what makes the untrusted search
cheap: a search can afford to explore many more boxes than the kernel will
ever see. And in the kernel FloatLib's scalar operations are currently two
to four times faster than \realfp's; \realfp{} is fast in kernel mode
because its certificates need few operations (derivative and Taylor
certificates instead of deep bisection), not because each operation is
cheaper. Adopting FloatLib-style kernel arithmetic is an obvious
improvement.

\paragraph{Shared goals with LeanCert and FloatLib.}
\Cref{tab:shared} lists 21 goals from a comparison with LeanCert
(bounds with and without touching, multivariate polynomials and
transcendental functions, integrals). All three tools were run in kernel
mode, one goal per file, with wall times including imports. \realfp{} and
LeanCert were measured on 23~November~2026 on the same machine (LeanCert
revision \lean{44b323a} on Lean~4.35; median of three; imports 3.4\,s and
3.5\,s respectively); FloatLib's \lean{interval} tactic was measured for this
paper (revision \lean{d24844f}, Lean~4.36, imports 2.7\,s, timeout
300\,s; median of three for successes; part of this run shared the machine
with an unrelated job). A rerun of the \realfp{} column on the day of the
FloatLib run, under the same load, again proved 20 of 21, with times
0.1--4\,s higher; the table keeps the earlier, uncontended run. ``---'' means the tactic did not prove the goal; this is a
statement about the automation on this exact statement, not about what can be
proved with more manual guidance.

\begin{table}[ht]
\centering\small
\caption{Shared goals, kernel mode: wall time in seconds including imports (top row). ``---'': not proved; n/s: integrals are not supported by FloatLib's \lean{interval} tactic.}
\label{tab:shared}
\begin{tabular}{@{}llrrr@{}}
\toprule
 & Goal & LeanCert & FloatLib & \realfp\\
\midrule
 & imports only & 3.48 & 2.7 & 3.43\\
\midrule
p01 & $\log 2<7/10$ & 3.44 & 2.66 & 3.34\\
p02 & $0.6931471<\log 2<0.6931472$ & 4.48 & --- & 3.51\\
p03 & $3.14159<\pi<3.1416$ & 3.62 & 2.58$^\dagger$ & 3.56\\
p04 & $2.7182\le e\le2.7183$ & 3.86 & 2.70$^\dagger$ & 3.44\\
p05 & $e^x\cos x\le3$ on $[0,1]$ & 3.97 & 2.60 & 4.22\\
p06 & $x(1-x)\le1/4+10^{-3}$ on $[0,1]$ & 3.75 & 2.77$^\dagger$ & 3.36\\
p07 & $\sin x\le x+10^{-3}$ on $[0,1]$ & ---$^\ddagger$ & 4.74$^\dagger$ & 4.23\\
p08 & $1+x+x^2/2\le e^x$ on $[0,1]$ (equality at 0) & --- & --- & 3.52\\
p09 & $\sin x\le x$ on $[0,2]$ (equality at 0) & --- & --- & 4.37\\
p10 & $\sin x\cos y\le1.01$ on $[0,1]^2$ & 24.78 & 2.52 & 4.82\\
p11 & $xy-x-y\ge-1.01$ on $[0,1]^2$ & --- & 8.29$^\dagger$ & 3.57\\
p12 & $xy+yz+zx\le3$ on $[0,1]^3$ & 4.49 & 2.41 & 3.35\\
p13 & $e^{xy}-x-y\ge-0.01$ on $[0,1]^2$ & --- & 5.01$^\dagger$ & 11.34\\
p14 & $0.74\le\int_0^1e^{-x^2}$ & 89.54 & n/s & 4.67\\
p15 & $0.7468\le\int_0^1e^{-x^2}$ & --- & n/s & 5.72\\
p16 & $\int_1^2 dx/x\le0.6932$ & --- & n/s & 4.77\\
p17 & $\int_0^1\arctan\le0.43882$ & --- & n/s & ---\\
p18 & $\int_0^1\arctan\le0.439$ & --- & n/s & 9.35\\
p19 & $\int_1^2 dx/x\le0.7$ & --- & n/s & 7.50\\
p20 & $x^2+y^2-xy\ge-0.01$ on $[-1,1]^2$ & --- & 3.06$^\dagger$ & 5.57\\
p21 & Rosenbrock $\ge-0.01$ on $[-2,2]\times[-1,3]$ & --- & 2.75 & 4.59\\
\midrule
 & proved & 9/21 & 12/21 & 20/21\\
\bottomrule
\end{tabular}
\smallskip
\parbox{0.95\linewidth}{\footnotesize $^\dagger$FloatLib with \lean{depth := 24} (default depth 8 failed), and for conjunctions after splitting with \lean{refine $\langle$?\_,~?\_$\rangle$}; all other FloatLib entries use the defaults. $^\ddagger$LeanCert's front door fails; its subdivision tactic on the rewritten goal $\sin x-x\le10^{-3}$ takes 63\,s.}
\end{table}

\paragraph{A downstream re-verification.}
To test \realfp{} on a real research proof, we re-verified the Coulomb
($s=1$) theorem of the Lean formalization accompanying Kryvonos, Liehr and
Taylor's work on energy minimization for eight points on the
sphere~\cite{em8} (release of 20~September~2026, Lean~4.33), replacing its
numerical part by \realfp{} certificates. The re-verification exists in a
\lean{native\_decide} version and in a kernel-only version whose main
theorems use only the three standard axioms, as do the original's. Clean
builds on an NVIDIA DGX Spark (20 cores; logged 23~November~2026, Mathlib
prebuilt) took 37\,min~35\,s wall time and 11.7~CPU-hours for the original,
6\,min~50\,s and 1.1~CPU-hours for the kernel-only \realfp{} version, and
55\,s and 254~CPU-seconds for the native version; the native \realfp{} version has 101
modules (7{,}560 lines) against the original's 1{,}195 modules. In the other direction, a
\realfp{} re-verification of the branch-and-bound in Tomar's four-body
proof~\cite{fourbody} builds in under two minutes on the same machine but
only in native mode: replaying its roughly $6\cdot10^{10}$ interval
operations in the kernel is estimated at hundreds of CPU-hours, whereas the
original formalization's purpose-built kernel certificates take 7.6
CPU-hours. Kernel mode is the default, but it is not free, and the
appropriate trust level depends on the size of the computation.


\section{Assurance}\label{sec:assurance}

\paragraph{No \lean{sorry}; axiom audit.} The default build of the Mathlib
layer includes an audit module that enumerates every declaration in the
\lean{RealFP} namespace and fails the build if any theorem depends on an
axiom other than \lean{propext}, \lean{Classical.choice} and
\lean{Quot.sound}; \lean{sorryAx} is never allowed. Only theorems in the
\lean{RealFP.Demo} and \lean{RealFP.Examples} namespaces may use
\lean{native\_decide}, and the audit counts them. At the time of writing the
full audit covers 2051 theorems; a script performs the same check for the
Mathlib-free core.

\paragraph{Differential testing.} The C runtime is a trust assumption only
in native mode, but it is tested heavily because the search depends on it
for completeness. Test executables compare every C kernel with the Lean
specification bit for bit on random inputs, special values (signed zeros,
subnormals, overflow, infinities, NaN) and hard cases; compare the portable
backend with AVX-512 and with an \lean{fesetround} reference, also for
aarch64 under QEMU; and check the transcendental bounds against
high-precision \texttt{mpmath} oracles. A regression test checks that kernel
and compiled evaluation of a checker agree and that a bound just past the
true enclosure is rejected in both modes.

\paragraph{Continuous integration.} Every push to the main branch builds
the core library and proofs, runs the tests and both axiom audits, and
compiles every tutorial snippet (about 11--16 minutes on GitHub-hosted
runners, using the portable backend); the slow example files and the full
audit run weekly.

\paragraph{Size.} The library consists of about 51{,}000 lines of Lean in
187 files and about 3{,}300 lines of C, plus the unmodified CORE-MATH
sources.

\section{Related work}\label{sec:related}

\paragraph{Arb / FLINT.} Arb~\cite{arb}, now part of FLINT~\cite{flint},
implements arbitrary-precision ball (midpoint--radius) arithmetic with an
exceptionally broad range of special functions, polynomials, power series,
matrices and rigorous integration, and it is very fast. It is the natural
tool for rigorous numerics outside a proof assistant, but its results are
not formal proofs: one trusts its C implementation. \realfp{} has far fewer
functions and only two precisions, but its results are Lean theorems; Arb
remains useful alongside it for exploration and for finding certificates.

\paragraph{Coq/Rocq.} The most mature formal counterpart is CoqInterval by
Melquiond and coauthors~\cite{melquiond2008,martindorelmelquiond2016},
which provides the \lean{interval} tactic for Coq (now Rocq~\cite{rocq})
with bisection, automatic differentiation, univariate Taylor models and
verified definite integrals~\cite{mahboubi2016}. It builds on the Flocq
formalization of floating-point arithmetic~\cite{flocq} and on the
Coquelicot real-analysis library~\cite{coquelicot}, and since primitive
binary64 floats were added to the Coq kernel~\cite{primfloats} it can
evaluate with hardware floats in reduction. \realfp{} follows the same
``compute, then check by reflection'' philosophy in Lean, where no such
library existed. Two differences: \realfp{} offers an explicit choice
between pure kernel evaluation and compiled evaluation per proof, and it
adds multivariate Krawczyk systems, exact positive definiteness and
touching-point certificates. CoqInterval's univariate Taylor-model
machinery is more developed.

\paragraph{Isabelle/HOL and HOL Light.} Isabelle's \lean{approximation}
method~\cite{hoelzl2009} proves real inequalities by interval arithmetic
with software floats. In HOL Light, Solovyev and Hales verified the
nonlinear inequalities of the Kepler conjecture with formal Taylor interval
approximations~\cite{solovyevhales2013,flyspeck}.

\paragraph{FloatLib.} FloatLib~\cite{floatlib} is a large Lean~4 library of
executable floating-point formats (IEEE binary and decimal, posits,
low-precision machine-learning formats) with proved rounding, status flags,
roundoff analysis, and kernel-checked automation: \lean{float\_decide}, and
an \lean{interval} tactic that proves real inequalities by certified
evaluation and subdivision, with quadratic forms for points where the margin
is zero. Its emphasis is breadth of formats and floating-point semantics;
its interval arithmetic runs on software floats in the kernel. At the
revision we examined it had no hardware directed-rounding path in its
proofs. \realfp{} adapts one routine from FloatLib (a bit-length
method used by the kernel fast path; MIT license).

\paragraph{LeanCert.} LeanCert~\cite{leancert} is a Lean~4 library for
certified numerics with rational, dyadic and affine interval backends,
automatic differentiation, global optimization, root existence and
Krawczyk systems, integration, finite sums, natural-number tails and a
Python SDK. It is broader than \realfp{} in goal shapes and tooling; its
default trust mode is \lean{native\_decide} with a kernel option. \realfp{}
is narrower but faster on the goals both handle (\cref{sec:bench}).

\paragraph{Other Lean work.} Irving's \lean{interval}
library~\cite{irvinginterval} provides conservative interval arithmetic in
Lean~4 on a software 64+64-bit floating-point type, with proofs of
conservativity.

\section{Limitations}\label{sec:limits}

We list the main limitations; a complete, current list is maintained in the
repository.
\begin{itemize}[nosep]
\item \emph{Shape of goals.} Domains are boxes with numeral endpoints, so
  $[0,\pi/2]$ must be replaced by a rational sub-interval or handled by a
  separate argument. Optimization certificates are first order, and argmin
  localization is one-dimensional. \lean{integral\_bound} handles
  one-dimensional integrals (a two-dimensional Taylor quadrature exists at
  the library level).
\item \emph{Dimension and dependency.} Like all interval methods, the cost
  of bisection grows quickly with the number of variables and with
  cancellation in the expression. Derivative, centred-form and touching
  certificates reduce the number of boxes greatly, but goals in more than a
  handful of variables need problem-specific structure.
\item \emph{Kernel cost.} A kernel-checked elementary function evaluation
  costs roughly 0.1\,s for $\log$ and 0.5\,s for $\sin$ or $\cos$, so tight
  trigonometric bounds needing many boxes can be slow or exceed heartbeat
  limits; a few such examples in the repository use native mode. Very large
  certificates (millions of operations) are beyond kernel checking and need
  \lean{trust := native}; see \cref{sec:bench}.
\item \emph{Touching points.} A bound attained \emph{exactly} at an
  irrational point, such as $(x^2-2)^2\ge0$ at $\sqrt2$, cannot be proved by
  interval evaluation at all, because every enclosure of $f$ near $\sqrt2$
  contains negative numbers; such goals need a symbolic step
  (\lean{nlinarith} suffices here). Degenerate minima such as $(x-1)^4$ are
  rejected, \lean{touch := auto} only finds interior non-degenerate minima
  (up to six variables), and touching certificates are binary64 only.
\item \emph{Precision 128.} Some functions ($\arctan$, $\tan$, $\arcsin$,
  $\arccos$, registered functions) are still evaluated with binary64
  accuracy inside the 128-bit mode, and the 128-bit Krawczyk preconditioner
  is a binary64 inverse.
\item \emph{Function coverage.} Compared with Arb, the set of functions is
  small; there are no special functions beyond $\operatorname{erf}$ and the
  registered examples, and no complex interval arithmetic.
\item \emph{Toolchain.} The library is pinned to Lean and Mathlib
  \lean{v4.36.0}. On x86-64 the default backend needs AVX-512F and gcc;
  other machines use the portable backend.
\item \emph{Native mode.} In native mode the C runtime is trusted on the
  strength of testing, not proof.
\end{itemize}


\section{Availability and conclusion}

\realfp{} is available at \url{https://github.com/schlafly1/realfp} under
the Apache-2.0 license. A project uses it by adding
\begin{lstlisting}
require realfp_mathlib from git
  "https://github.com/schlafly1/realfp" @ "v0.1.1" / "RealFPMathlib"
\end{lstlisting}
to its lakefile and using the same Lean and Mathlib versions
(\lean{v4.36.0}). The repository contains a tutorial in which every snippet
is compiled, a reference manual, the soundness notes summarized in
\cref{cor:trust}, the benchmark scripts and logs behind \cref{sec:bench},
and a list of known limitations.

The design principle is simple: specifications in Lean, speed in C,
checking in the kernel. Floating-point arithmetic is precisely specified by
IEEE~754, so it can be modelled exactly in a proof assistant; directed
rounding makes it rigorous; and the search/check split lets the fast,
unverified part do the hard work of finding certificates while only a small
verified checker has to be trusted. For a mathematician with a numerical
lemma inside a formal proof, the result is a tactic that is usually one
line, usually takes seconds, and produces a theorem with the same axioms as
the rest of Mathlib.

\bibliographystyle{plain}
\bibliography{refs}

\end{document}
